\documentclass[10pt]{article}

\usepackage[margin=0.6in]{geometry}
\usepackage{enumitem}
\usepackage{titlesec}
\usepackage{xcolor}
\usepackage{hyperref}
\definecolor{linkblue}{HTML}{0B5FA5}
\hypersetup{colorlinks=true, urlcolor=linkblue, linkcolor=linkblue, citecolor=linkblue}
\usepackage{parskip}

\titleformat{\section}{\large\bfseries}{}{0em}{}[{\vspace{-5pt}\rule{\linewidth}{0.6pt}}]
\titlespacing{\section}{0pt}{6pt}{3pt}
\setlist[itemize,1]{leftmargin=1.2em, itemsep=1pt, topsep=1pt}
\setlist[itemize,2]{leftmargin=0.75em, labelsep=0.28em, itemsep=0.6pt, topsep=1pt,
                    label=\raisebox{0.2ex}{\rule{0.25em}{0.25em}}}
\pagestyle{empty}
\setlength{\parskip}{2pt}
\newcommand{\me}{\textbf{M.\ Islam}}

\begin{document}

%========================= HEADER =========================
\begin{center}
    {\LARGE\bfseries Mainul Islam}\\[3pt]
    \href{mailto:mainulislam@iut-dhaka.edu}{mainulislam@iut-dhaka.edu} \;\textbar\;
    \href{https://scholar.google.com/citations?user=mHfGHugAAAAJ\&hl=en}{Google Scholar} \;\textbar\;
    \href{https://www.linkedin.com/in/mainul-islam-254071205/}{LinkedIn} \;\textbar\;
    \href{https://mainul-islam-07.github.io/mainulislam.github.io/}{Portfolio} \;\textbar\;
    \href{https://www.youtube.com/@mainulislamefan}{YouTube}
\end{center}

%========================= EDUCATION =========================
\section*{Education}
\textbf{B.Sc.\ in Electrical and Electronic Engineering} \hfill Jan 2020 -- Jun 2024\\
\href{https://www.iutoic-dhaka.edu/}{Islamic University of Technology (IUT), Gazipur, Bangladesh}
\hfill CGPA: 3.87 / 4.00

\smallskip
\textit{Thesis -} 
\begin{itemize}
    \item MonoSIM, an open-source SIL platform developed for testing autonomous control algorithms on Ackermann vehicles.
    % \item Contributed to the monocular lane-detection pipeline: camera calibration, bird's-eye-view transform, and sliding-window polynomial lane fitting
    \item  The system was designed using open source tools, creating an environment with a monocular camera vision system to capture stimuli from it with minimal computational overhead through a sliding window based lane detection method.
    \item The system is validated with PID vs. MPC controllers, showing PID had lower tracking error while MPC gave smoother, constraint-aware control.
\end{itemize}

% Developed an open-source simulation framework for autonomous vehicles
% using camera-based lane detection, evaluated predictive and classical control strategies for
% path tracking (MPC, PID). Published and presented.

%========================= PROFESSIONAL EXPERIENCE =========================
\section*{Professional Experience}

\textbf{Research and Development Engineer} ---
\href{https://cyberneticsbd.com/}{Cybernetics Hi-Tech Solutions Pvt.\ Ltd.}, Dhaka
\hfill Oct 2024 -- Present
\begin{itemize}
    \item Contributed to the development of robotic and embedded systems from initial requirements
    through to delivery, covering hardware, software, and operator interfaces, for applications like industrial-automation etc.
    \item Design and implemented autonomous navigation stack for ground vehicles, including
    positioning, route planning, motion control, and the coordination of multi-vehicle fleets.
    \item Developed real-time embedded system for distributed motor control and sensor systems
    across a range of industrial communication standards.
    % \item Define the safety and fault-tolerance architecture of critical platforms, including
    % interlocks, automated fault recovery, and redundant communication links.
    \item Designed custom control electronics from schematic through fabrication, testing, and
    commissioning for demanding field environments.
    % \item Conduct system integration, field trials, and client handover, and prepare the
    % supporting technical documentation.
\end{itemize}

\textbf{Lecturer (Part-Time)} ---
\href{https://www.aust.edu/}{Ahsanullah University of Science and Technology}, Dhaka
\hfill Jul 2025 -- Jun 2026
\begin{itemize}
    \item Conducted C++ Programming and Digital Signal Processing courses.
    \item Delivered lectures, prepared laboratory manuals, and conducted student assessment.
\end{itemize}

%========================= PUBLICATIONS =========================
\section*{Publications}

\textbf{Peer-Reviewed Journal Articles}
\begin{itemize}
    \item A.\ U.\ R.\ Adib, \me, M.\ S.\ Abid, R.\ Ahshan, ``A deep learning based framework for
    solar panel segmentation and fault classification enhanced with explainable AI,''
    \textit{Solar Energy}, vol.\ 302, 2025. (Impact Factor: 6.6)
    DOI: \href{https://doi.org/10.1016/j.solener.2025.114058}{10.1016/j.solener.2025.114058}
    \item S.\ Rahman, A.\ Banik, M.\ K.\ Hossain, \me, A.\ Al-Hysam, I.\ Ahmed, ``Topological survey
    of zeta power converters from classical designs to emerging architectures,''
    \textit{Energy Conversion and Management: X}, art.\ no.\ 102222, 2026. (Impact Factor: 8.8)
    DOI: \href{https://doi.org/10.1016/j.ecmx.2026.102222}{10.1016/j.ecmx.2026.102222}
\end{itemize}

\textbf{Peer-Reviewed Conference Papers}
\begin{itemize}
    \item S.\ Rahman, N.\ Hasin, \me, M.\ Z.\ A.\ Rony, G.\ Sarowar, ``MonoSIM: An open-source SIL
    framework for Ackermann vehicular systems with monocular vision,'' \textit{IEEE 12th Intl.\
    Conf.\ on Automation, Robotics and Applications (ICARA)}, 2026.
    DOI: \href{https://doi.org/10.1109/ICARA69401.2026.11480327}{10.1109/ICARA69401.2026.11480327}
    \item N.\ Hasin, M.\ A.\ R.\ Nishat, \me, K.\ S.\ A.\ Hasan, A.\ Newaz, ``A robust deep learning
    framework for Bangla license plate recognition using YOLO and vision-language OCR,''
    \textit{IEEE Intl.\ Conf.\ on AI and Data Analytics (ICAD)}, 2026.
    DOI: \href{https://doi.org/10.1109/ICAD69378.2026.11608728}{10.1109/ICAD69378.2026.11608728}
    \item H.\ Haque, \me, M.\ A.\ I.\ Alvy, M.\ S.\ Islam, M.\ A.\ Razzak, ``Long-term energy demand
    analysis using machine learning algorithms: A case study in Bangladesh,'' \textit{IEEE ICECIE},
    2024.
    DOI: \href{https://doi.org/10.1109/ICECIE63774.2024.10815643}{10.1109/ICECIE63774.2024.10815643}
    \item M.\ S.\ Tanvir, F.\ A.\ Tonmoy, M.\ M.\ Islam, \me, M.\ Hasan, ``Optimizing master--slave
    broadcast and unicast communication in multi-node STM32 networks using UART and RS485
    protocols,'' \textit{IEEE ICCIT}, 2024.
    DOI: \href{https://doi.org/10.1109/ICCIT64611.2024.11022430}{10.1109/ICCIT64611.2024.11022430}
   
\end{itemize}
\newpage
\textbf{Manuscripts Under Review}
\begin{itemize}
    
    \item M.\ M.\ Islam, \me, M.\ S.\ Tanvir, ``Speed synchronization of a dual-BLDC tracked
    differential-drive vehicle using a GA-tuned multi-rate cross-coupled controller.''
\end{itemize}


%========================= PROJECTS =========================
\section*{Projects}
\begin{itemize}[label={}, leftmargin=0pt, itemsep=5pt]

    \item \textbf{\href{https://youtu.be/fLiwTvl2xmI}{Man-Transportable EOD Robot
    (Jontro Soinik 2.0)}}
    \begin{itemize}
        % \item One of the principal developers of an explosive ordnance disposal vehicle with a seven-degree-of-freedom manipulator.
        \item Developed the motion-control library governing eleven servo drives distributed over
        two independent CAN buses, implementing the CANopen (CiA-402) standard so that each joint
        can be commanded by position, speed, or torque.
        \item Contributed to a closed-form solution to the arm's inverse kinematics. Verified
        against the robot's geometric model and integrated with MoveIt\,2 for collision checking.
        \item Wrote the STM32 firmware carrying the operator's control signal over an Ethernet
        link via optical fiber.
        % \item Designed the layered fail-safe architecture: a fault on one bus is isolated without
        % disabling the whole robot, the platform holds a defined safe state on error, and the
        % disruptor disarms the instant the operator link drops.
        \item Developed an app interface, presenting low-latency
        video feeds, a live telemetry dashboard, and a three-dimensional on-screen model mirroring
        the robot's actual posture.
    \end{itemize}

    \item \textbf{\href{https://youtu.be/TkmOu6lCBTw?si=8KOdXk-vVMPzP4vm}{Autonomous Guided Vehicle ---
    CYBERFLEET}}
    \begin{itemize}
        \item Designed the complete ROS\,2 system integration for a three-vehicle
        warehouse robot fleet, covering how each vehicle establishes its position, plans a route,
        follows it, and coordinates with the others.
        \item Implemented an extended Kalman filter combining continuous wheel-encoder motion with
        intermittent camera sightings of fixed floor markers.
        \item Built the marker-based localisation pipeline: camera calibration, pose estimation
        from each observed marker.
        \item Contributed to the route planning that treats time as a planning dimension, preventing both collisions and deadlocks.
        % \item Engineered the central fleet coordinator, which allocates jobs, runs an independent
        % state machine for each vehicle, and enforces exclusive access to shared lift and scanner
        % stations.
        \item Implemented LiDAR-based obstacle detection.
    \end{itemize}

    \item \textbf{\href{https://www.youtube.com/watch?v=MLrD0SuG4Q0}{Autonomous Target-Tracking
    Drone}}
    \begin{itemize}
        \item Contributed to the implementation of live onboard target detection \& tracking, running a Pi\,5 with an AI accelerator. 
        \item Created the aerial training
        dataset and tested different tracking methods.
        % \item Optimised thrust distribution and flight profile to extend mission endurance, and
        % validated flight stability and mission execution in both simulation and field trials.
    \end{itemize}



    \item \textbf{\href{https://github.com/Mainul-Islam-07/Tollbooth_Automation}{Toll-Booth Automation}}
    \begin{itemize}
        \item Developed the lane-controller firmware on an STM32 microcontroller, structured as a
        table-driven state machine.
        % so that the sequence from vehicle detection to payment
        % verification to barrier release is explicit, auditable, and readily adapted per site.
        \item Integrated inductive loop coils, an infrared exit sensor, and licence-plate
        recognition to detect vehicles.
        % confirm payment, and raise security-breach alerts
        % automatically.
        % \item Hardened the communication links against the electrically noisy roadside
        % environment, with automatic detection and recovery from corrupted serial data and an
        % interrupt-driven Ethernet server that keeps response timing predictable under load.
        \item Designed three custom circuit boards with optical isolation, surge and reverse-polarity protection.
        % so that electrical faults on the roadside cannot reach the controller; included the drive
        % stages for boom barriers, sirens, and indicator lighting.
        % \item Structured the firmware so that every lane and every site is served from a single
        % code base through build-time configuration, avoiding separate software versions per unit,
        % with manual operator override at each step and a defined safe-reset path.
    \end{itemize}

    % \item \textbf{\href{https://youtu.be/bgiEhOPq0wE}{STM32 Based Ethernet Communication}}
    % \begin{itemize}
    %     \item Built a reliable SBUS communication system over both optical fiber and wireless links using WIZnet modules and STM32 (Bluepill).
    %     \item Designed a two-stage interrupt pipeline (ISR flags events, main loop handles them) to keep the ISR non-blocking.

    % \end{itemize}

    % \item \textbf{\href{https://www.youtube.com/watch?v=3r5VhLD1Ths}{Jontro Soinik 1.0} and Smart
    % Soinik 1.0}
    % \begin{itemize}
    %     % \item Developed embedded firmware for real-time actuation and sensor feedback loops on the
    %     % second and third generation of Bangladesh's indigenous explosive ordnance disposal
    %     % vehicles.
    %     \item Contributed to an addressable serial-bus architecture in which every control node carries
    %     its own address.
    %     \item Contributed to sensor integration
    %     and to power system design through to circuit-board fabrication.
    % \end{itemize}

    % \item \textbf{Surveillance Drone (SD15-N):} Quantified how aerodynamic drag, motor thrust
    % against current draw, altitude, and battery capacity each limit flight endurance, and used the
    % resulting model to optimise battery selection and weight distribution; confirmed the
    % improvement through field testing.

    % \item \textbf{\href{https://www.youtube.com/watch?v=bSHUT5XJSNQ}{Offline Hardware Cryptocurrency
    % Wallet}}
    % \begin{itemize}
    %     \item Developed a fully crypto transaction-signing device on which private keys are
    %     generated and stored locally and never reach a networked computer, with all data crossing the
    %     gap as scanned QR codes.
    %     Implemented the device state machine, touchscreen interface,
    %     standards-based (BIP-39) recovery-phrase generation with word-by-word confirmation, and
    %     multi-page QR output for signed transactions.
    % \end{itemize}
    % \item \textbf{Fixed-Wing VTOL UAV:} Performed the analytical sizing of thrust, lift and drag,
    % aspect ratio, wing area, and stall speed; configured the flight controller across four
    % aerodynamic control surfaces; and contributed to manufacture, assembly, and flight testing of a
    % two-metre prototype.

    % \item \textbf{Undergraduate projects:} Designed an educational 8-bit computer (arithmetic unit,
    % control unit, registers) executing a shortened instruction cycle; developed a MATLAB
    % battery-management simulation as well as tested hardware that isolates the load whenever any cell leaves its safe
    % window, scalable across two to six-cell packs; and wrote microcontroller assembly firmware measuring signal frequency and duty cycle
    % with threshold alarms and generating a signal at the sum of two input frequencies.

    % \item \textbf{IUT Mars Rover ---
    % \href{https://www.facebook.com/projectaltair}{Project Altair} and
    % \href{https://www.facebook.com/iutavijatrik}{Team Avijatrik}:} Served as electrical system
    % designer across two competition teams --- designed the wheel-drive board and several other PCBs
    % and sub-circuits, and handled wiring, hardware integration, and debugging.

\end{itemize}

\vspace{-3pt}
%========================= TECHNICAL SKILLS =========================
\section*{Technical Skills}
\vspace{-3pt}
\textbf{Programming:} Python, C/C++, Embedded C, MATLAB. \\
\textbf{Robotics and Control:} ROS\,2, sensor fusion, vision-based perception. \\
\textbf{Embedded Systems:} STM32 platforms, CAN, UART, Modbus.\\
\textbf{Machine Learning:} PyTorch, TensorFlow, scikit-learn, object detection, tracking,
explainable AI. \\
\textbf{Design Tools:} EAGLE, KiCAD, easy eda.
%========================= HONORS =========================
\section*{Honours and Achievements}
\begin{itemize}
    \item \textbf{International Rover Challenge 2024:} 6th globally and 1st in Bangladesh; awarded
    ABEX Best Science Team. 
    % (Project Altair).
    \item \textbf{European Rover Challenge 2023:} 16th globally (Project Altair).
    \item \textbf{\href{https://www.mathworks.com/matlabcentral/profile/authors/25804044}{MATLAB
    Cody}:} Global rank 117, with 914 problems solved and 26 badges.
\end{itemize}

%========================= REFERENCES =========================
\section*{References}
\textbf{\href{https://scholar.google.com/citations?hl=en\&user=3waAWaYAAAAJ\&view_op=list_works\&sortby=pubdate}{Prof.\ Dr.\ Golam Sarowar}} --- B.Sc.\ Thesis Supervisor \\
Professor, Department of EEE, Islamic University of Technology (IUT) \\ 
Email: asim@iut-dhaka.edu

\end{document}